\chapter{ABSTRACT}
\indent\indent In the era of widespread digitalization, digital documents such as invoices, legal contracts, certificates, and identity cards serve as the cornerstone of financial and administrative transactions. However, the advancement of modern image editing tools has enabled increasingly sophisticated text tampering, severely threatening automated verification workflows. Classical digital forensic techniques often struggle to distinguish subtle local alterations, whereas modern deep learning models are typically single-task oriented (confined to either binary classification or pixel segmentation) and suffer from excessive parameter counts, hindering deployment on resource-constrained platforms.

To address these limitations, this study proposes \textbf{LightweightDTD}, a lightweight multi-task deep learning architecture comprising only approximately \textbf{4.15 million parameters (4,155,796 parameters)}. The proposed framework integrates two complementary information streams: (1) a \textbf{Spatial Branch} powered by an EfficientNet-B0 backbone to extract multi-level visual features; and (2) a \textbf{Frequency Branch} employing an orthonormal \textbf{8x8 BlockDCT (Block Discrete Cosine Transform)} to capture fine-grained JPEG compression inconsistencies and edge interpolation artifacts. These multi-domain features are harmoniously fused through a concurrent \textbf{scSE (Spatial and Channel Squeeze-and-Excitation)} module and subsequently fed into a \textbf{LiteMID (Lightweight Multi-view Iterative Decoder)} to aggregate contextual representations across multiple receptive fields ($H/4$ to $H/32$). Operating under a multi-task learning paradigm, LightweightDTD concurrently delivers two outputs: a pixel-level tampering segmentation mask ($224 \times 224$) and multi-label classification over six microscopic forensic artifact categories.

Extensive experiments on the standardized \textbf{Minh-TFR TestSet} (comprising 2,000 independent test images) demonstrate the superior performance of the proposed architecture: in pixel-level segmentation, LightweightDTD achieves a \textbf{Pixel Accuracy of 99.79\%}, a \textbf{Micro Precision of 90.09\%}, a \textbf{Micro Recall of 77.93\%}, a \textbf{Micro IoU of 71.78\%}, and a \textbf{Macro Dice score of 62.18\%}. In multi-label tampering artifact classification, the model attains a \textbf{Micro F1 of 77.49\%}, a \textbf{Macro F1 of 58.60\%}, a \textbf{mAP (PR-AUC) of 61.02\%}, and a \textbf{Top-3 Class Hit Rate of 98.45\%}. These empirical results validate that LightweightDTD provides accurate localization and multi-attribute forensic characterization of tampered document text while preserving exceptional computational efficiency, facilitating practical integration into edge document verification systems.\vspace{0.3cm}

\textbf{Keywords:} Document Image Forensics, Tampered Text Detection, LightweightDTD, EfficientNet-B0, BlockDCT, scSE Attention, LiteMID, Multi-task Learning, Multi-label F1.
